\documentclass[a4paper,14pt]{extreport}

\usepackage{graphicx}

% Пакеты для русского языка и кодировки
\usepackage[T2A]{fontenc}
\usepackage[utf8]{inputenc}
\usepackage[russian]{babel}

% Пакеты для математики и теорем
\usepackage{amsmath, amssymb, amsthm}

% Пакеты для графики и гиперссылок
\usepackage{hyperref}

% Настройка полей
\usepackage{geometry}
\geometry{top=2cm, bottom=2cm, left=2.5cm, right=1.5cm}

% Для оформления кода и таблиц
\usepackage{listings}
\usepackage{longtable,booktabs,calc}
\usepackage[table]{xcolor}

% Настройка межстрочного интервала
\usepackage{setspace}
\onehalfspacing

% Отключает нумерацию для part, chapter, section и т.д.
\setcounter{secnumdepth}{-3}

\hypersetup{
    colorlinks=true,
    linkcolor=black,
    urlcolor=blue
}

\usepackage{tikz}
% Пересечения проводов без точки не являются соединениями.
\newcommand{\krCommonScheme}{%
\begin{center}
\begin{tikzpicture}[line width=.85pt, every node/.style={font=\small}]
\node[draw, minimum width=1.35cm, minimum height=.8cm] (or) at (2,2.45) {ИЛИ};
\node[draw, minimum width=2.2cm, minimum height=.9cm, align=center] (xor) at (2,1.4) {исключ.\\ИЛИ};
\node[draw, minimum width=1.15cm, minimum height=.7cm] (invb) at (2,.15) {НЕ};
\node[draw, minimum width=1cm, minimum height=.8cm] (and1) at (4.4,1.9) {И};
\node[draw, minimum width=1cm, minimum height=.8cm] (and2) at (6.2,1.05) {И};
\node[draw, minimum width=1cm, minimum height=.7cm] (inv) at (7.7,1.05) {НЕ};
\draw (0,3) node[left] {$A$} -- (.45,3) -- (.45,1.6) -- (.9,1.6);
\draw (.45,2.65) -- (1.325,2.65);
\fill (.45,2.65) circle (1.5pt);
\draw (0,2) node[left] {$B$} -- (.6,2) -- (.6,2.25) -- (.6,.15) -- (1.425,.15);
\draw (.6,2.25) -- (1.325,2.25);
\fill (.6,2) circle (1.5pt);
\draw (0,1) node[left] {$C$} -- (.8,1) -- (.8,1.2) -- (.9,1.2);
\draw (2.675,2.45) -- (3.4,2.45) -- (3.4,2.1) -- (3.9,2.1);
\draw (3.1,1.4) -- (3.35,1.4) -- (3.35,1.7) -- (3.9,1.7);
\draw (4.9,1.9) -- (5.35,1.9) -- (5.35,1.25) -- (5.7,1.25);
\draw (2.575,.15) -- (5.25,.15) -- (5.25,.85) -- (5.7,.85);
\draw (6.7,1.05) -- (7.2,1.05);
\draw (8.2,1.05) -- (8.7,1.05) node[right] {$F$};
\end{tikzpicture}
\end{center}}

\newcommand{\krShiftScheme}{%
\begin{center}
\begin{tikzpicture}[line width=.85pt, every node/.style={font=\small}]
\node[draw, minimum width=2.2cm, minimum height=.9cm, align=center] (xor) at (2,2.5) {исключ.\\ИЛИ};
\node[draw, minimum width=1cm, minimum height=.8cm] (and1) at (4.1,1.8) {И};
\node[draw, minimum width=1.1cm, minimum height=.8cm] (or) at (6,1.2) {ИЛИ};
\node[draw, minimum width=1cm, minimum height=.7cm] (inv) at (7.7,1.2) {НЕ};
\node[draw, minimum width=1cm, minimum height=.8cm] (and2) at (9.5,1) {И};
\draw (0,3) node[left] {$A$} -- (.45,3) -- (.45,-.25) -- (8.85,-.25) -- (8.85,.8) -- (9,.8);
\fill (.45,2.7) circle (1.5pt);
\draw (.45,2.7) -- (.9,2.7);
\draw (0,2) node[left] {$B$} -- (.75,2) -- (.75,2.3) -- (.9,2.3);
\draw (0,1) node[left] {$C$} -- (3.1,1) -- (3.1,1.6) -- (3.6,1.6);
\draw (3.1,2.5) -- (3.4,2.5) -- (3.4,2) -- (3.6,2);
\draw (4.6,1.8) -- (5.05,1.8) -- (5.05,1.4) -- (5.45,1.4);
\draw (0,.1) node[left] {$D$} -- (4.9,.1) -- (4.9,1) -- (5.45,1);
\draw (6.55,1.2) -- (7.2,1.2);
\draw (8.2,1.2) -- (8.6,1.2) -- (8.6,1.2) -- (9,1.2);
\draw (10,1) -- (10.5,1) node[right] {$F$};
\end{tikzpicture}
\end{center}}


\begin{document}

\begin{center}
    \textbf{Контрольная работа. Вариант 1}
\end{center}

\begin{enumerate}
    \item Представьте число $3260.21_8$ в шестнадцатеричной системе счисления.

    \item Переведите число $2128_{10}$ в шестнадцатеричную систему счисления.

    \item Переведите число $0.65625_{10}$ в восьмеричную систему счисления.

    \item Вычислите сумму $123.44_5 + 332.11_5$. Ответ запишите в пятеричной системе счисления.

    \item Даны числа $A=342_{10}$ и $B=-430_{10}$. Длина разрядной сетки --- 16 двоичных разрядов. Представьте числа в дополнительном коде и выполните сложение $A+B$. Полученный 16-битный дополнительный код результата запишите восьмеричными цифрами.

    \item Даны числа $A=-1011_8$ и $B=56_{10}$. В основной памяти ЭВМ эти числа представлены в формате с фиксированной точкой в дополнительном коде. Длина формата --- 16 двоичных разрядов; все разряды, кроме знакового, отведены под целую часть числа. Выполните операцию $A-B$ и запишите полученный 16-битный дополнительный код результата шестнадцатеричными цифрами.

    \item Представьте число $123.344_{10}$ в формате одинарной точности IEEE 754 (\texttt{float32}). Укажите знаковый бит, поле порядка и поле дробной части мантиссы. Используйте округление к ближайшему представимому значению; при равенстве расстояний --- к значению с чётным младшим битом значащей части.

    \item Составьте таблицу истинности данной логической схемы. Упрощать и минимизировать выражение не требуется. Пересечение проводов без точки не означает соединения.
    \krCommonScheme
\end{enumerate}

\end{document}
